%%% Vyplňte prosím základní údaje o závěrečné práci (odstraňte \xxx{...}).
%%% Automaticky se pak vloží na všechna místa, kde jsou potřeba.

% Druh práce:
%	"bc" pro bakalářskou
%	"mgr" pro diplomovou
%	"phd" pro disertační
%	"rig" pro rigorozní
\def\ThesisType{bc}

% Název práce v jazyce práce (přesně podle zadání)
\def\ThesisTitle{Bitcoinová peněženka pro pokročilé uživatele}

% Název práce v angličtině
\def\ThesisTitleEN{Bitcoin Wallet for Advanced Users}

% Jméno autora (vy)
\def\ThesisAuthor{Jakub Dvořák}

% Rok odevzdání
\def\YearSubmitted{2026}

% Název katedry nebo ústavu, kde byla práce oficiálně zadána
% (dle Organizační struktury MFF UK:
% https://www.mff.cuni.cz/cs/fakulta/organizacni-struktura,
% případně plný název pracoviště mimo MFF)
\def\Department{Katedra softwarového inženýrství}
\def\DepartmentEN{Department of Software Engineering}

% Jedná se o katedru (department) nebo o ústav (institute)?
\def\DeptType{Katedra}
\def\DeptTypeEN{Department}

% Vedoucí práce: Jméno a příjmení s~tituly
\def\Supervisor{RNDr. Filip Zavoral, Ph.D.}

% Pracoviště vedoucího (opět dle Organizační struktury MFF)
\def\SupervisorsDepartment{Katedra softwarového inženýrství}
\def\SupervisorsDepartmentEN{Department of Software Engineering}

% Studijní program (kromě rigorozních prací)
\def\StudyProgramme{Informatika}

% Nepovinné poděkování (vedoucímu práce, konzultantovi, tomu, kdo
% vám nosil pizzu a vařil čaj apod.)
\def\Dedication{%
Rád bych poděkoval vedoucímu práce RNDr. Filipu Zavoralovi, Ph.D.
za cenné rady a~připomínky při vedení této práce.
}

% Abstrakt (doporučený rozsah cca 80-200 slov; nejedná se o zadání práce)
\def\Abstract{%
Tato práce se zabývá návrhem a~implementací mobilní aplikace pro platformu
Android, která umožňuje pokročilou správu Bitcoinových prostředků. Aplikace
podporuje multisignature transakce ve schématu M-of-N s~využitím standardů
BIP-48 a~BIP-67, manuální výběr nepoužitých výstupů (coin control) a~integraci
s~hardwarovou peněženkou Trezor prostřednictvím protokolu Trezor Connect.
Bezpečnostní model zajišťuje, že server nikdy nedisponuje privátními klíči --
veškeré podepisování probíhá výhradně na~hardwarovém zařízení. Pro distribuci
částečně podepsaných transakcí mezi cosignery je využit formát PSBT dle
standardu BIP-174. Backend je realizován jako soustava mikroslužeb v~jazyce
Kotlin s~frameworkem Ktor, frontend jako nativní Android aplikace s~Jetpack
Compose. Funkčnost aplikace byla ověřena na~testnetové síti včetně úspěšného
dokončení 2-of-3 multisig transakce.
}

% Anglická verze abstraktu
\def\AbstractEN{%
This thesis presents the design and implementation of a~mobile application
for the Android platform that enables advanced Bitcoin asset management.
The application supports multisignature transactions in an M-of-N scheme
using the BIP-48 and BIP-67 standards, manual selection of unspent transaction
outputs (coin control), and integration with the Trezor hardware wallet via
the Trezor Connect protocol. The security model ensures that the server never
holds private keys -- all transaction signing is performed exclusively on the
hardware device. The PSBT format defined by BIP-174 is used for distributing
partially signed transactions among cosigners. The backend is implemented as
a~set of microservices in Kotlin with the Ktor framework, while the frontend
is a~native Android application built with Jetpack Compose. The application
was verified on the Bitcoin testnet, including a~successful completion of
a~2-of-3 multisig transaction.
}

% 3 až 5 klíčových slov (doporučeno) oddělených \sep
% Hodí se pro nalezení práce podle tématu.
\def\ThesisKeywords{%
Bitcoin\sep multisignature\sep PSBT\sep hardwarová peněženka\sep Android
}

\def\ThesisKeywordsEN{
Bitcoin\sep multisignature\sep PSBT\sep hardware wallet\sep Android
}

% Pokud některá z položek metadat obsahuje TeXové řídící sekvence, je potřeba
% dodat i verzi v obyčejném textu, která se objeví v metadatech formátu XMP
% zabudovaných do výstupního souboru PDF. Pokud si nejste jistí, prohlédněte si
% vygenerovaný soubor thesis.xmpdata.
\def\ThesisAuthorXMP{\ThesisAuthor}
\def\ThesisTitleXMP{Bitcoinova penezenka pro pokrocile uzivatele}
\def\ThesisKeywordsXMP{Bitcoin, multisignature, PSBT, hardwarova penezenka, Android}
\def\AbstractXMP{\Abstract}

% Máte-li dlouhý abstrakt a nechceme se mu vejít na informační stranu,
% můžete použít toto nastavení ke zmenšení písma informační strany.
% (Uvažte nicméně zkrácení abstraktu, to je často lepší.)
\def\InfoPageFont{}
%\def\InfoPageFont{\small}  % odkomentujte pro zmenšení písma
